\chapter{Мгновенные платежи в~мире: PIX, UPI, FedNow, SEPA Instant}\label{ch:instant-payments-world}

СБП из предыдущей главы~-- одна из национальных систем мгновенных
безотзывных переводов между банковскими счетами в~режиме 24/7/365.
Такие же системы работают в~Бразилии (Pix), Индии (UPI), США (FedNow)
и~еврозоне (SEPA Instant Credit Transfer). Глава разбирает их
параллельно по одинаковым параметрам: оператор расчёта, способ
инициации платежа (есть ли QR в~правилах схемы), роль небанковских
участников и~встроенный в~схему антифрод.

\section{Pix (Бразилия)}\label{sec:pix-system}

Pix~-- платёжная схема, оператором которой выступает сам центральный
банк Бразилии (BCB). Pix запущен 16 ноября 2020~года и~с~момента
старта получил обязательное покрытие: участие обязательно для всех
финансовых институтов, у~которых более чем 500~тыс.\ активных счетов
клиентов~\cite{pix,pix-participants}.

Расчёт идёт через систему мгновенных платежей SPI~--- выделенную
RTGS-инфраструктуру под управлением BCB. В~режиме RTGS (real-time
gross settlement, валовой расчёт в~реальном времени) каждый перевод
проходит индивидуально и~окончательно, без сведения в~нетто-позицию
и~без расчётных окон~--- в~отличие от карточного клиринга, где сотни
тысяч операций между парой участников сводятся в~одну нетто-позицию
в~конце дня (см.~\ref{sec:recon-net-vs-gross}). Каждый участник держит
в~SPI отдельный расчётный счёт; финальность наступает в~момент
дебета счёта плательщика и~кредита счёта получателя~\cite{pix-spi}.
Частного клирингового слоя между участниками нет.

\subsection*{Ключи и~DICT}\label{subsec:pix-keys-dict}

Реквизит, по которому отправляется Pix, не привязан к~номеру счёта
напрямую. У плательщика на руках~--- ключ Pix: номер мобильного
телефона, e-mail, CPF (идентификационный номер физлица), CNPJ
(идентификатор юрлица) или случайный 32-символьный идентификатор.
Директория DICT под управлением BCB отображает ключ в~реквизиты
счёта и~участника. Перед отправкой банк плательщика обращается
к~DICT, получает реквизиты получателя и~формирует расчётную
инструкцию в~SPI.

\subsection*{Инициация: QR-код и~платёжная строка}\label{subsec:pix-init}

BCB определяет форматы инициации Pix в~отдельном
регламенте~\cite{pix-manual-iniciacao}. Статический QR-код многоразовый,
содержит ключ получателя и~опционально сумму. Динамический QR
одноразовый, генерируется под конкретную транзакцию и содержит ссылку
на платёжную сессию у~PSP получателя; кассовое ПО получает
обратный вызов по факту исполнения. Помимо сканирования, тот же набор
полей передаётся как текстовая платёжная строка, которую копируют
в~буфер обмена и~вставляют в~банковское приложение~\cite{pix-faq}.
Этот режим закрывает сценарий, когда и~PSP-приложение, и~источник
кода открыты на одном устройстве и~сканировать QR камерой невозможно.

Формат QR-кода Pix задаёт сам BCB: EMVCo-совместимый BR Code определён
на уровне самой схемы. С~продуктовой точки зрения это означает, что
приложения разных банков понимают один и~тот же код, и~подписи плательщика
на разных каналах ведут к~одному и~тому же сообщению в~SPI.

\subsection*{Тайминги и~SLA}\label{subsec:pix-timings}

BCB регулирует тайминги Pix отдельным регламентом, периодически
пересматривая разрешённые длительности стадий~\cite{pix-manual-tempos}.
Документ задаёт максимальные интервалы на отдельные этапы~--- от запроса
к~DICT до завершения расчёта~--- и~служит SLA в~правовом смысле.
Сообщения между участниками и~SPI идут по защищённому каналу через
сертификаты ICP-Brasil; требования к~подписи и~каналу прописаны
в~отдельном регламенте по безопасности~\cite{pix-manual-seguranca}.

\subsection*{Возвраты и~MED}\label{subsec:pix-med}

Pix~--- push-перевод с~немедленной финальностью, поэтому обычный
возврат идёт как новая, обратная по направлению, операция Pix. Сценарий
\enquote{я отправил не туда} в~Pix отдельно урегулирован: получатель
обязан рассмотреть запрос на возврат, но согласие на него остаётся
за получателем, если иное не установлено по закону или суду.

Для случаев мошенничества BCB ввёл специальный механизм возврата
MED: когда у~банка-отправителя есть фактические основания считать
операцию мошеннической, банк инициирует через инфраструктуру Pix
запрос на блокировку и~возврат суммы у~получателя. Регламент
исполнения Pix описывает MED как отдельный поток внутри
схемы~\cite{pix-manual-fluxos}. По продуктовой семантике MED работает
примерно так же, как chargeback в~карте.

\section{UPI (Индия)}\label{sec:upi-system}

UPI запущен NPCI в~апреле 2016~года и~работает поверх IMPS~---
более раннего сервиса 24/7-переводов. Регулирует Резервный банк
Индии (RBI) по закону о~платёжных и~расчётных системах
2007~года~\cite{upi}. Технически UPI~--- набор открытых API, которыми
пользуются банки и~небанковские приложения, плюс центральный
коммутатор NPCI, маршрутизирующий запросы и~участвующий в~клиринге.

\subsection*{VPA и~четырёхсторонняя модель}\label{subsec:upi-fourparty}

Реквизит платежа в~UPI~--- виртуальный платёжный адрес VPA, формата
\texttt{name@psp} (например, \texttt{abhishek@okaxis}). VPA отображается
на конкретный банковский счёт без раскрытия номера счёта плательщику.
Аналогично DICT в~Pix, VPA снимает с~пользователя задачу запоминать
банковские реквизиты.

В~отличие от карточной четырёхсторонней схемы, в~UPI участвуют четыре
типа сущностей: банк-плательщика, банк-получателя, PSP-bank (банк-спонсор,
который держит контракт с~NPCI и~отвечает за TPAP-приложение) и~TPAP
(небанковское приложение с~клиентским интерфейсом). PhonePe, Google Pay, Paytm~--- крупнейшие
TPAP; PSP-bank над ними даёт регуляторное прикрытие. NPCI в~центре
выступает коммутатором: маршрутизирует запросы, идентифицирует
участников через VPA, ведёт клиринг и~расчёт.

\subsection*{Антифрод и~рыночные ограничения}\label{subsec:upi-tpap-cap}

Чтобы не допустить концентрации, NPCI установил лимит: один TPAP
не может превышать 30~\% месячного объёма UPI-транзакций. Процедура мониторинга
закреплена отдельным циркуляром~\cite{npci-upi-tpap-cap}: при
превышении NPCI обязан вводить ограничительные меры на новых
пользователях TPAP до возврата доли в~рамки порога. Это рыночное
ограничение поверх технического протокола; ничего подобного в~правилах
СБП, Pix, FedNow или SEPA Instant нет.

В~2025~году NPCI ужесточил операционные требования к~использованию
UPI-API~\cite{npci-upi-oc-215,npci-upi-circulars}. Действующий
циркуляр ограничивает частоту вызова вспомогательных API: проверка
баланса~--- не более 50 раз в~сутки на пользователя, запрос списка
счетов~--- не более 25, инструкции AutoPay выполняются в~трёх
определённых окнах вместо непрерывного потока. Обязательны ежегодный
независимый аудит и~квартальная самоаттестация участников. Эти
правила живут в~регуляторном слое поверх протокола UPI и~сдерживают
DDoS-нагрузку на банковские бэкенды со стороны массовых TPAP.

\subsection*{Сценарии и~отличие от Pix}\label{subsec:upi-vs-pix}

UPI поддерживает push (\enquote{я плачу}) и~collect (\enquote{запросите
с~меня}) с~подтверждением плательщика. QR-инфраструктура в~UPI
стандартизирована (Bharat QR / UPI QR), но конкуренция между приложениями
смещает массовый сценарий в~сторону набора VPA и~deep-link. В~Pix
фронт замкнут на банковские приложения, и~поток зависит от BCB.
В~UPI центральный коммутатор принадлежит NPCI, а~клиентский слой
собирается из конкурирующих TPAP-приложений.

\section{FedNow (США)}\label{sec:fednow-system}

FedNow запущен 20 июля 2023~года и~эксплуатируется Федеральными
резервными банками~\cite{fednow-about,fednow-resources}. До его
появления розничный американский ландшафт делили Fedwire (межбанковский
RTGS для крупных сумм, не 24/7 до недавнего расширения), ACH
(батчевые межбанковские переводы, T+1) и~частный RTP The Clearing
House (мгновенный, с~2017~года, но не центробанковский). FedNow
закрывает этот разрыв центробанковским сервисом, доступным всем
держателям мастер-счёта в~ФРС.

\subsection*{Правовая основа}\label{subsec:fednow-legal}

FedNow регулируется Subpart~C 12~CFR Part~210~\cite{reg-j-subpart-c}
и~Operating Circular~8~\cite{fednow-oc8}. Платёжная инструкция
в~FedNow определена как сообщение, переданное через сервис; правила
Subpart~C вытесняют непротиворечащие положения Article~4A UCC. FedNow подпадает
прямо под федеральный закон; рамки ответственности участников
жёстче, чем у~частного RTP, который работает по правилам собственной
системы.

Operating Procedures~\cite{fednow-op-procedures} задают операционные
требования: 24-часовой операционный день, режим 7/365, никаких
плановых окон отключения. Для участников это означает, что отдельного
закрытия дня в~операционном смысле нет: суточная сверка отделена
от расчётного цикла.

\subsection*{ISO 20022 как родной формат}\label{subsec:fednow-iso20022}

В~отличие от UPI или СБП, у~которых протоколы свои, FedNow построен
поверх ISO 20022 с~момента запуска. Банк-плательщика передаёт
\texttt{pacs.008} в~FedNow; FedNow дебетует счёт банка-плательщика
в~ФРС, кредитует счёт банка-получателя; банк-получателя зачисляет
средства клиенту. Подтверждение возвращается через \texttt{pacs.002},
уведомления клиентам идут через
\texttt{camt.054}~\cite{fednow-iso20022}. Руководство по~имплементации
и~песочница для тестирования сообщений размещены на платформе
MyStandards SWIFT~\cite{fednow-tech-overview}.

\subsection*{Расчёт и~ликвидность}\label{subsec:fednow-settlement}

Расчёт в~FedNow валовой, в~деньгах центрального банка, через
мастер-счёт участника в~ФРС. Сетевой расчётный риск здесь отсутствует:
расчётчиком выступает сам ФРС, и~дефолт расчётчика невозможен
(в~отличие от частного RTP). У~участника возникает другая задача~---
поддерживать на мастер-счёте достаточно ликвидности круглосуточно,
включая ночи и~выходные. Параллельно с~запуском FedNow ФРС расширил
доступ к~ликвидности (продлённые часы дисконтного окна и~сервисы
внутридневного кредитования), чтобы участники не выпадали из системы
в~часы низкой ликвидности.

\subsection*{Чего у~FedNow нет}\label{subsec:fednow-gaps}

Своего QR-кода у~FedNow нет: схема ничего не говорит про инициацию~---
ни о~ссылке для оплаты, ни о~запросе на оплату с~деталями назначения
платежа, ни о~форматах QR. Это оставлено рынку~--- банковским
ядрам, финтех-приложениям, маркетплейсам. На практике участники
строят пользовательский слой поверх FedNow самостоятельно или через
надстройку RTP, где такие сценарии существуют.

\section{SEPA Instant Credit Transfer (еврозона)}\label{sec:sepa-instant-system}

SCT Inst~--- мгновенный кредитовый перевод в~евро, определённый
правилами Европейского платёжного совета (EPC). Целевое сквозное
время~--- 10~секунд, 24/7/365. В~еврозоне работают две расчётные
платформы: центробанковская TIPS (оператор Eurosystem)~\cite{ecb-tips}
и~частная RT1 (оператор EBA Clearing). Участник PSP подключается
к~одной из них; между ними настроена взаимная достижимость.

\subsection*{IPR: обязательность через регулирование}\label{subsec:sepa-ipr}

До 2024~года участие в~SCT Inst было добровольным; за восемь лет
после запуска (ноябрь 2017) покрытие так и~не стало полным. Это
изменилось с~принятием Regulation (EU) 2024/886~--- регламента
о~мгновенных платежах (IPR)~\cite{eu-2024-886,ecb-ipr}. Регламент
опубликован в~Официальном журнале ЕС 19~марта 2024~года, вступил
в~силу 8~апреля 2024~года и~ввёл сроки обязательной реализации:

\begin{itemize}
  \item 9~января 2025~--- PSP в~еврозоне обязаны принимать
        SCT Inst;
  \item 9~октября 2025~--- PSP в~еврозоне обязаны отправлять
        SCT Inst и~предоставлять услугу проверки получателя (VOP);
  \item 9~января 2027~--- PSP вне еврозоны обязаны принимать;
  \item 9~июля 2027~--- PSP вне еврозоны обязаны отправлять;
  \item 9~июня 2028~--- финальные обязательства, в~том числе по тарифной
        нейтральности и~доступу небанковских PSP к~платёжным системам.
\end{itemize}

IPR убирает разрешённую ранее наценку за SCT Inst поверх обычного
SCT и~открывает прямой доступ к~платёжным системам для лицензированных
небанковских PSP.

\subsection*{Правила SCT Inst 2025}\label{subsec:sepa-rulebook}

EPC опубликовал свод правил SCT Inst 2025 v1.0 в~ноябре 2024~года;
вступил в~силу 5~октября 2025~года вместе с~основной волной
обязательств IPR~\cite{epc-sct-inst-rulebook-2025}. Документ согласован
с~IPR: выровнены поля, тайминги и~обработка отказов. Историческое
ограничение по сумме на одну операцию (100~000~EUR) IPR снимает~---
лимит определяется правилами расчётной системы и~PSP.

\subsection*{Проверка получателя (VOP)}\label{subsec:sepa-vop}

Параллельно с~основным сводом EPC ввёл отдельную схему проверки
получателя (VOP). Перед инициированием SCT/SCT Inst PSP плательщика
спрашивает PSP получателя, соответствует ли указанное имя получателя
реквизитам IBAN. Ответ~--- \emph{совпадение}, \emph{близкое совпадение}
(с~предложенным именем), \emph{нет совпадения} или \emph{проверка
невозможна}~--- передаётся плательщику до подтверждения
платежа~\cite{epc-vop}. Целевое время отклика VOP не более одной
секунды, верхняя планка~--- три секунды. Центральный каталог PSP,
который ведёт EPC, маршрутизирует VOP-запрос на нужную точку приёма.

VOP~--- регуляторно обязательная мера антифрода: IPR требует
от PSP её реализации с~9~октября 2025~года. Для интеграторов это
означает, что клиентский интерфейс A2A-платежа в~еврозоне обязан
явно показывать результат VOP перед подтверждением, а~PSP плательщика
отвечает за корректное поведение, если плательщик настаивает
на отправке при ответе \enquote{нет совпадения}.

\section{Сравнение систем мгновенных платежей}\label{sec:instant-comparison}

В~таблице~\ref{tab:instant-comparison} системы мгновенных платежей
сведены по одинаковым параметрам, включая СБП из главы~\ref{ch:sbp}.

\begin{table}[tbp]
  \begingroup
  \scriptsize
  \setlength{\tabcolsep}{3pt}
  \renewcommand{\arraystretch}{1.2}
  \begin{tabularx}{\textwidth}{
      @{}>{\RaggedRight\arraybackslash\bfseries}p{0.16\textwidth}
      >{\RaggedRight\arraybackslash}p{0.155\textwidth}
      >{\RaggedRight\arraybackslash}p{0.155\textwidth}
      >{\RaggedRight\arraybackslash}p{0.155\textwidth}
      >{\RaggedRight\arraybackslash}p{0.155\textwidth}
      >{\RaggedRight\arraybackslash}p{0.155\textwidth}@{}
    }
    \toprule
    \rowcolor{Panel}
    & СБП & Pix & UPI & FedNow & SCT Inst \\
    \midrule
    \rowcolor{Soft}
    Оператор & НСПК (Банк России) & BCB & NPCI (под RBI) & Федеральные резервные банки & TIPS (Eurosystem) и~RT1 (EBA Clearing) \\
    Расчёт & ЦБ через ОПКЦ & BCB через SPI & коммутатор NPCI и~банки-расчётчики & мастер-счёт в~ФРС & TIPS~--- ЦБ-деньги; RT1~--- ЦБ-расчёт по итогам \\
    \rowcolor{Soft}
    Запуск & 2019 & ноябрь 2020 & апрель 2016 & 20 июля 2023 & ноябрь 2017 \\
    Протокол & ОПКЦ OpenAPI & проприетарный поверх ICP-Brasil & UPI API NPCI & ISO 20022 (pacs.008) & ISO 20022 (pacs.008) \\
    \rowcolor{Soft}
    QR в~правилах схемы & да (НСПК) & да (BR Code, BCB) & да (UPI QR, NPCI) & нет & нет \\
    Обязательность & добровольная для банков, обязательная для крупных в~ряде сценариев & обязательная с~500~тыс.\ счетов & регулируется RBI и~NPCI & добровольная & обязательная по IPR с~2025 \\
    \rowcolor{Soft}
    Антифрод схемы & период охлаждения & MED & лимит суммы 30~\%, лимит частоты API & нет на уровне схемы & VOP (IBAN + имя) \\
    Возврат & новая операция и~процедура у~банка & новая операция, MED для фрода & новая операция, диспут через PSP & новая операция & новая операция, плюс отзыв по правилам схемы \\
    \bottomrule
  \end{tabularx}%
  \endgroup
  \caption{Сравнение систем мгновенных платежей по одинаковым параметрам.}\label{tab:instant-comparison}
\end{table}

\section{Что общее и что разное}\label{sec:instant-lessons}

ISO 20022 распространён, но не универсален: у~UPI собственный API
NPCI, у~Pix~--- собственный протокол поверх ICP-Brasil, в~СБП~---
ОПКЦ OpenAPI; FedNow и~SCT Inst построены прямо на ISO 20022.

\textbf{Кто оператор расчёта}. В~Pix и~СБП оператором выступает
центральный банк страны с~собственной инфраструктурой расчёта.
FedNow эксплуатируется самими Федеральными резервными банками. SCT Inst работает
на двух платформах (центробанковской TIPS и~частной RT1), между
которыми настроена взаимная достижимость. В~UPI центральный коммутатор
принадлежит NPCI, расчёт между банками идёт под надзором RBI. Расчёт в~деньгах центрального банка убирает сетевой расчётный
риск, но требует от регулятора нести операционную ответственность
за круглосуточный сервис.

\textbf{Где живёт уровень инициации}. Pix и~UPI определяют форматы
QR-кода и~эквивалентов (платёжная строка в~Pix, deep-link в~UPI)
в~правилах самой схемы. СБП~--- аналогично, через НСПК. FedNow и~SCT Inst оставляют
инициацию рынку: схема видит только финальную инструкцию pacs.008.
Когда клиентский слой встроен в~правила схемы, опыт между PSP
стандартизирован; когда он вынесен за пределы схемы, остаётся место
для конкуренции, но усложняется совместимость.

\textbf{Роль небанковских участников}. UPI ввёл явный класс TPAP
и~ограничил его 30~\% объёма~\cite{npci-upi-tpap-cap}, чтобы клиентский
слой оставался конкурентным. IPR в~еврозоне открывает прямой доступ
к~платёжным системам лицензированным небанковским PSP. В~Pix, СБП
и~FedNow небанковские приложения работают через банка-спонсора,
и~отдельного класса участника схемы у~них нет.

\textbf{Антифрод как часть схемы}. VOP в~SCT Inst работает \emph{до}
инициации (предотвращение), MED в~Pix~--- \emph{после} (восстановление).
В~СБП есть период охлаждения на крупные операции при подозрении. В~UPI ограничения работают через регуляторные
циркуляры~\cite{npci-upi-oc-215}: лимит частоты вызова API, лимит
суммы для одного TPAP, обязательные аудиты. У~FedNow антифрод на уровне схемы отсутствует;
ответственность переложена на банка-отправителя и~банка-получателя.

Различия в~правилах возврата, лимитах, антифроде и~уровне
инициации значимы для бэк-офиса: сверка, поддержка, диспут-процедуры,
требования к~клиентскому интерфейсу и~к~хранению доказательств
собираются под конкретную систему отдельно.
